\section{Extensions}
\label{sec:extensions}

Beyond the core, \PurePy is designed to grow through extensions that add expressiveness while preserving
purity. This section sketches two, each drawn from \Fluid: partial application and piecewise function
definitions. We give an example of each. \todo{Formally define or defer to future work}

\subsection{Partial application}
\label{sec:extension-partial-application}

A function applied to fewer arguments than it declares returns a function of the remaining arguments:

\begin{python}
def f(x, y):
    return x + y

f(5)(6)  # 11
\end{python}

\noindent The syntax is valid Python, but Python raises a \pyinline{TypeError} at runtime, because
\pyinline{f(5)} is missing an argument. For partial application to add meaning rather than change it, core
\PurePy must reject such unsaturated calls, requiring every call to supply the declared arity. On that basis,
\pyinline{f(5)} denotes the function that adds \pyinline{5} to its argument, so \pyinline{f(5)(6)} evaluates to
\pyinline{11}. The extension is therefore semantic rather than syntactic: it gives a meaning to programs that
the core rejects.

\subsection{Piecewise function definitions}
\label{sec:extension-piecewise}

A function may be defined by several clauses that pattern-match on its arguments, in place of a single
definition with an explicit \pyinline{match}:

\begin{python}
def length(Nil):
    return 0
def length(Cons(x, xs)):
    return 1 + length(xs)
\end{python}

\noindent Python has no way to define a function by pattern-matching on its arguments: \pyinline{def
length(Cons(x, xs))} is a syntax error, since a parameter cannot be a pattern. Core \PurePy, a subset of
Python, has no such form either. The extension adds this: the clauses together define \pyinline{length}, dispatching on
the constructor of the argument. Because Python gives the pattern-clause form no meaning, it is a
syntactic extension that needs no special notation to distinguish it from an ordinary \pyinline{def}.
